\subsection{Experimental Causal Discovery}

\subsubsection{Causal Effects in the Real World}

\begin{Eg}[Does the thermometer cause the sun to rise?]
    \label{therm-sun}
    Consider an old type of thermometer ($T$) with a needle that can - for simplicity of arguments - either point to higher temperatures (up) or to lower temperatures (down).
    We also consider the state of the sun ($S$), which can either be up ($u$) or down ($d$).
    We then observe that $T$ correlates with $S$. For simplicity, we assume a one-to-one relationship:
    \[\begin{array}{c|c} 
        T & S \\ 
        \hline
        u & u \\
        d & d
    \end{array}\]
    The conditional distribution $P(S|T)$ then looks like this:
    \begin{align*}
        P(S=u|T=u) & = 1, \\
        P(S=u|T=d) & = 0, \\
        P(S=d|T=u) & = 0, \\
        P(S=d|T=d) & = 1.
    \end{align*}
    If we are cold we are now tempted to try changing the needle in the thermometer in order to make the sun rise and warm us up.\\
What is wrong with our analysis?
\end{Eg}

\begin{Disc}
    The example \ref{therm-sun} makes clear that there is a difference between:
    \begin{enumerate}
        \item \emph{Observing} the movement of the the thermometer needle $T$ and the sun $S$, using 1 \emph{observational} data set, 
            leading to an estimate of $P(S|T)$.
        \item \emph{Interacting} with the thermometer needle $T$ and getting the sun's response $S$, using probably many \emph{interventional} data sets, leading to estimates for $P(S|\doit(T))$:
                \begin{align*}
                    P(S=u|\doit(T=u)) & = 0.5, \\
                    P(S=u|\doit(T=d)) & = 0.5, \\
                    P(S=d|\doit(T=u)) & = 0.5, \\
                    P(S=d|\doit(T=d)) & = 0.5.
    \end{align*}
    \end{enumerate}
\end{Disc}

\begin{Def}[Causal effect---real world definition]
    \label{def-causality-real}
    We say that a variable $X$ has \emph{causal effect} onto another variable $Y$ if
    when \emph{forcing} $X$ to take different values $x_0$, $x_1$ then also the distribution of $Y$ changes, in symbols:
    \[ \exists x_0,x_1 \in \Xcal:\qquad P(Y|\doit(X=x_0)) \neq P(Y|\doit(X=x_1)). \]
\end{Def}

\begin{Rem}
    \begin{enumerate}
        \item 
    Again, note that example \ref{therm-sun} shows that the condition in definition \ref{def-causality-real} is different from:
    \[ \exists x_0,x_1 \in \Xcal:\qquad P(Y|X=x_0) \neq P(Y|X=x_1), \]
    which just uses the conditional distribution instead of the interventional distribution.
        \item Also note, that the '$\doit$-operators' are not operators on the observational distribution $P(X,Y)$ or $P(Y|X)$, etc., or on the corresponding observational data sets. They reflect actions/interventions in the real world leading to different distributions and corresponding data sets. 
        \item There are usually many possible intervention values and targets one can think of leading to many different interventional distributions and data sets.
        \item One can consider the observational distribution as a special case of an interventional distribution (where we intervene by doing nothing).
    \end{enumerate}
\end{Rem}


\subsubsection{Randomized Controlled Trials (RCT)} % - The gold standard for experimental causal discovery }


\begin{Prn}[Randomized Controlled Trial (RCT)]\label{prn:rct}
 Assume we want to know if variable $X$ has causal influence on outcome variable $Y$, i.e.\ we want to estimate the deviation between:
$ P(Y|\doit(X=x_0))$ and  $P(Y|\doit(X=x_1))$. For this we have test subjects $w_1,\dots,w_N$.
A \emph{Randomized Controlled Trial} then follows the following steps:
\begin{enumerate}
    \item Split the population of test subjects into 2 groups ('test group' $C_1$ vs. 'control group' $C_0$) by random lot (or fair coin flips).
    \item Give every test subject $w_n \in C_1$ from 'test group' the treatment $x_1$ and the ones  $w_m \in C_0$ from 'control group' the control treatment $x_0$.
    \item Measure the outcome $y_n$ for each test subject $w_n$ and estimate the deviation: \[D:=d(P(Y|\doit(X=x_0)),P(Y|\doit(X=x_1))).\]
    \item Do a statistical test if the deviation $D$ is significantly different from $0$. 
    \item If it is significantly different we can conclude a \emph{causal effect} of $X$ onto $Y$, otherwise not.
\end{enumerate}
\end{Prn}


\begin{Eg}
    Example application of randomized controlled trials are:
    \begin{enumerate}
        \item drug or vaccine testing,
        \item advertisement placement,
        \item evaluating public policies, etc.
        \item A. Banerjee, E. Duflo, M. Kremer got the Nobel Prize in Economics 2019 for using RCTs in poverty research, e.g.\: improving school attendance and performance in poor areas via giving different towns different incentives (e.g.\ text books vs. deworming medicine vs. control groups).
    \end{enumerate}
\end{Eg}

\begin{Disc}
    \begin{enumerate}
        \item An RCT is an 'interventional study' (in contrast to just 'observational study') since we can control the treatment and 'force' it onto the test subjects.
        \item Randomized Controlled Trials are considered the gold standard for experimental causal discovery.
        \item To further avoid biases one usually insists on double/triple blind RCT studies, 
            i.e.\ noone directly involved in the study knows who got which treatment 
            (e.g.\ neither the doctor, the experimenter, the patient, etc.). 
        \item Often RCTs cannot be done for ethical reasons (e.g.\ "smoking causes cancer" research).
        \item Sometimes RCTs require too many resources to be feasible.
    \end{enumerate}
\end{Disc}


\begin{Exc}
    Go online, find news like "drinking wine every day is good for your health" or "chewing gum causes diabetes", etc., 
    look up the original research paper and check:
    \begin{enumerate}
        \item if they did interventional studies (like RCT) or just observational studies,
        \item if and/how they did doubly/triply blind RCTs (in case they did one),
        \item otherwise, if and/or how they ruled out other correlation inducing scenarios,
        \item what bias could have possibly introduced through the data collecting process,
        \item how big the data set was, what assumptions were made, what statistical methods were used, etc.,
        \item what other 'stories' you could come up with in order to explain the data. Be creative, create 5 stories!
    \end{enumerate}
    Write your findings down and talk to others about it.
\end{Exc}
